\documentclass[11pt]{article}
\usepackage[a4paper,margin=21mm]{geometry}
\usepackage{fontspec}
\setmainfont{Latin Modern Roman}
\newfontfamily\CJKfont{Noto Serif CJK SC}
\XeTeXlinebreaklocale "zh"
\XeTeXlinebreakskip=0pt plus 1pt
\usepackage{amsmath,amssymb,amsthm,amscd,graphicx,color}
\usepackage{xurl}
\usepackage[unicode,hidelinks]{hyperref}
\allowdisplaybreaks
\setlength\emergencystretch{3em}
\numberwithin{equation}{section}
\newtheorem{Theorem}{Theorem}[section]
\newtheorem*{Theorem*}{Theorem}
\newtheorem{Corollary}[Theorem]{Corollary}
\newtheorem{Lemma}[Theorem]{Lemma}
\newtheorem{Proposition}[Theorem]{Proposition}
 { \theoremstyle{definition}
\newtheorem{Definition}[Theorem]{Definition}
\newtheorem{Note}[Theorem]{Note}
\newtheorem{Example}[Theorem]{Example}
\newtheorem{Remark}[Theorem]{Remark} }

\newcommand{\field}[1]{\mathbb{#1}}
\newcommand{\N}{\field{N}}
\newcommand{\Z}{\field{Z}}
\newcommand{\R}{\field{R}}
\newcommand{\C}{\field{C}}
\newcommand{\K}{\textbf{K}}

\begin{document}
\begin{center}\Large Bilateral Bailey lattice with independently varying parameters\end{center}
\noindent\textbf{Stable ID:} 560\quad\textbf{Author:} Jehanne Dousse; Frédéric Jouhet; Isaac Konan\par
\noindent\textbf{Work:} Bilateral Bailey Lattices and Andrews–Gordon Type Identities\par
\noindent\textbf{Source:} \url{https://sigma-journal.com/2025/032/}\par
\noindent\textbf{Version:} Official SIGMA 21 (2025) PDF and native TeX; acquired 2026-09-30; exact bytes pinned by SHA256\par
\noindent\textbf{Rights:} CC BY 4.0. Authors retain copyright. \url{https://creativecommons.org/licenses/by/4.0/}. Reuse and typesetting adaptation; original language and mathematical content preserved.\par
\noindent\textbf{Difficulty (prior selection preserved):} {\CJKfont H2: The independently varying parameter lattice requires deriving a symmetric-polynomial/q-binomial recurrence and proving coefficient cancellation in two layers before the rank induction closes. Repeated application of the one-parameter Bailey lemma alone does not establish the explicit finite alpha transform; adjacent specialized identities are grouped with this construction rather than counted separately.}\par
\medskip\noindent\textbf{Editorial boundary note (not author text):} Complete original Theorem 1.16, all bilateral/q-product/convergence/finite-sum conventions, full symmetric-polynomial coefficient recurrence proof and final dimension induction. The old range contained only the proof; the actual theorem at source268–298 is now included. The two displayed q-binomial product lemmas are retained as routine algebraic prerequisites with the author’s explicit disclosure: their proofs are straightforward verifications and are therefore omitted here. The nontrivial coefficient decompositions and cancellation-to-recurrence argument are fully written by the author and preserved; no assistant verification or missing-proof generation is added. Easier one-parameter subcases and applications are not additional problem counts.\par
\medskip\hrule\medskip\noindent\textbf{Author text begins below.}\par
\setcounter{section}{0}
% Original source lines 69--83
\section{Introduction and statement of results}\label{sec:intro}

A classical approach to obtain and prove $q$-series identities is the Bailey lemma, originally found by Bailey~\cite{Ba}, and whose iterative strength was later highlighted by Andrews~\cite{A1, A2, AAR} through the so-called Bailey chain.
Fix complex numbers $a$ and $q$. Recall~\cite{Ba} that a Bailey pair $((\alpha_n)_{n\geq0}, (\beta_n)_{n\geq0})$ ($(\alpha_n, \beta_n)$ for short) relative to $a$ is a pair of sequences satisfying
\begin{equation}\label{bp}
\beta_n=\sum_{j=0}^n\frac{\alpha_j}{(q)_{n-j}(aq)_{n+j}}\qquad\forall n\in\mathbb{N}.
\end{equation}
Here and throughout the paper, we use standard $q$-series notations which can be found in~\cite{GR}
\begin{equation*}
(a)_\infty = (a;q)_\infty:=\prod_{j\geq 0}\bigl(1-aq^j\bigr)\qquad\mbox{and}\qquad(a)_k = (a;q)_k:=\frac{(a;q)_\infty}{\bigl(aq^k;q\bigr)_\infty},
\end{equation*}
where $k \in \Z$, and
$
(a_1,\ldots,a_m)_k:=(a_1)_k\cdots(a_m)_k$,
where $k$ is an integer or infinity, and as usual~${|q|<1}$ to ensure convergence of infinite products.

\setcounter{Theorem}{4}\setcounter{equation}{3}
% Original source lines 133--245
\begin{Lemma}[McLaughlin]
\label{lem:Mac}
If $(\alpha_n, \beta_n)$ is a Bailey pair relative to $a$, then $(\alpha'_n, \beta'_n)$ is a Bailey pair relative to $a/q$, where
\begin{equation*}
\alpha'_0=\alpha_0, \qquad \alpha'_n=(1-a)\left(\frac{\alpha_n}{1-aq^{2n}}-\frac{aq^{2n-2}\alpha_{n-1}}{1-aq^{2n-2}}\right), \qquad \beta'_n=\beta_n.
\end{equation*}
\end{Lemma}

In this paper, we will show, among other things, that this lemma and the Bailey lattice can be extended to bilateral versions.


As noted in \cite{BMS} and~\cite{J}, it is possible to define for all $n\in\mathbb{Z}$ a \emph{bilateral Bailey pair} $(\alpha_n, \beta_n)$ relative to $a$ by the relation
\begin{equation}\label{bbp}
\beta_n=\sum_{j\leq n}\frac{\alpha_j}{(q)_{n-j}(aq)_{n+j}}\qquad\forall n\in\mathbb{Z}.
\end{equation}

\begin{Remark}\label{rk:bilatinversion}
The relation~\eqref{bp} defining classical (unilateral) Bailey pairs is a special instance of the above relation defining bilateral ones, as choosing $\alpha_n=0$ for negative integers $n$ in~\eqref{bbp} implies $\beta_n=0$ for $n$ negative. Actually, the converse is also true, as the classical Bailey inversion holds for bilateral Bailey pairs: $(\alpha_n, \beta_n)$ is a bilateral Bailey pair relative to $a$ if and only if
\begin{equation}\label{bilatinversion}
\alpha_n=\frac{1-aq^{2n}}{1-a}\sum_{j\leq n}\frac{(a)_{n+j}}{(q)_{n-j}}(-1)^{n-j}q^{\binom{n-j}{2}}\beta_j\qquad\forall n\in\mathbb{Z}.
\end{equation}
\textit{Thus, from all our results in this paper, one can deduce the corresponding unilateral results by setting $\alpha_n=0$ $($or equivalently $\beta_n=0)$ for all $n<0$.}
\end{Remark}

In \cite{BMS}, the Bailey lemma is extended in the following way.
\begin{Theorem}[bilateral Bailey lemma]\label{thm:bilatbaileylemma}
If $(\alpha_n, \beta_n)$ is a bilateral Bailey pair relative to $a$, then
so is $(\alpha'_n, \beta'_n)$, where
\begin{equation*}\alpha'_n=\frac{(\rho,\sigma)_n(aq/\rho\sigma)^n}{
(aq/\rho,aq/\sigma)_n}\alpha_n\qquad\mbox{and}\qquad\beta'_n=\sum_{j\leq n}\frac{(\rho,\sigma)_j(aq/\rho\sigma)_{n-j}(aq/\rho\sigma)^j}{(q)_{n-j}(aq/\rho,aq/\sigma)_n}\beta_j,
\end{equation*}
subject to convergence conditions on the sequences $\alpha_n$ and $\beta_n$, which make the relevant infinite series absolutely convergent.
\end{Theorem}

Our first result is an extension of the Bailey lattice to the bilateral case.

\begin{Theorem}[bilateral Bailey lattice]\label{thm:bilatbaileylattice}
If $(\alpha_n, \beta_n)$ is a bilateral Bailey pair relative to $a$, then~$(\alpha'_n, \beta'_n)$ is a bilateral Bailey pair relative to $a/q$, where
\begin{equation*}\alpha'_n=\frac{(\rho,\sigma)_n(a/\rho\sigma)^n}{(a/\rho,a/\sigma)_n}(1-a)\left(\frac{\alpha_n}{1-aq^{2n}}-\frac{aq^{2n-2}\alpha_{n-1}}{1-aq^{2n-2}}\right),
\end{equation*}
and
\begin{equation*}\beta'_n=\sum_{j\leq n}{(\rho,\sigma)_j(a/\rho\sigma)_{n-j}(a/\rho\sigma)^j\over (q)_{n-j}(a/\rho,a/\sigma)_n} \beta_j,
\end{equation*}
subject to convergence conditions on the sequences $\alpha_n$ and $\beta_n$, which make the relevant infinite series absolutely convergent.
\end{Theorem}

As mentioned above, several proofs have been given for the (unilateral) Bailey lattice. On the other hand, simpler proofs were given to prove the Andrews--Gordon identities without using the Bailey lattice. Here we give a very simple proof of our bilateral Bailey lattice, which when considering Bailey pairs such that $\alpha_n=0$ for $n<0$ reduces to the unilateral Bailey lattice. Hence we provide in particular a very simple proof of the classical Bailey lattice.

The key in our proof is the following simple lemma, which generalises McLaughlin's unilateral Lemma \ref{lem:Mac} and transforms bilateral Bailey pairs relative to $a$ into bilateral Bailey pairs relative to $a/q$.
\begin{Lemma}[key Lemma~1]\label{lem:key1}
If $(\alpha_n, \beta_n)$ is a bilateral Bailey pair relative to $a$, then $(\alpha'_n, \beta'_n)$ is a bilateral Bailey pair relative to $a/q$, where
\begin{equation}\label{eq:baileylattice1}
\alpha'_n=(1-a)\left(\frac{\alpha_n}{1-aq^{2n}}-\frac{aq^{2n-2}\alpha_{n-1}}{1-aq^{2n-2}}\right), \qquad \beta'_n=\beta_n,
\end{equation}
subject to convergence conditions on the sequences $\alpha_n$ and $\beta_n$, which make the relevant infinite series absolutely convergent.
\end{Lemma}
While it is not customary to do so, we give the proof in this introduction to show that it is just a few lines long and only requires the definition of bilateral Bailey pairs and elementary sum manipulations which are similar to the ones for the unilateral version.
\begin{proof}[Proof of Lemma~\ref{lem:key1}]
For all $n\in\mathbb{Z}$, we have
\begin{align*}
\sum_{j\leq n} \frac{\alpha'_j}{(q)_{n-j}(a)_{n+j}}&= \sum_{j\leq n} \frac{(1-a)}{(q)_{n-j}(a)_{n+j}}\left(\frac{\alpha_j}{1-aq^{2j}}-\frac{aq^{2j-2}\alpha_{j-1}}{1-aq^{2j-2}}\right)\\
&= \sum_{j\leq n} \frac{(1-a)\alpha_j}{(q)_{n-j}(a)_{n+j}\bigl(1-aq^{2j}\bigr)}-\sum_{j\leq n} \frac{(1-a)\bigl(1-q^{n-j}\bigr)aq^{2j}\alpha_{j}}{(q)_{n-j}(a)_{n+j+1}\bigl(1-aq^{2j}\bigr)}\\
&=\sum_{j\leq n} \frac{(1-a)\alpha_j}{(q)_{n-j}(a)_{n+j+1}\bigl(1-aq^{2j}\bigr)}\bigl(\bigl(1-aq^{n+j}\bigr)-aq^{2j}\bigl(1-q^{n-j}\bigr)\bigr)\\
&=\sum_{j\leq n} \frac{\alpha_j}{(q)_{n-j}(aq)_{n+j}}=\beta_n=\beta'_n,
\end{align*}
which is the desired result by~\eqref{bbp}.
\end{proof}
With this lemma, we can give an extremely simple proof of the bilateral Bailey lattice.
\begin{proof}[Proof of the bilateral Bailey lattice]
Start from a bilateral Bailey pair $(\alpha_n, \beta_n)$ relative to~$a$. Then apply Lemma~\ref{lem:key1} to obtain a bilateral Bailey pair $\bigl(\tilde{\alpha}_n, \tilde{\beta}_n\bigr)$ relative to $a/q$ and satisfying~\eqref{eq:baileylattice1}. Applying the bilateral Bailey lemma (see Theorem~\ref{thm:bilatbaileylemma}) to \smash{$\bigl(\tilde{\alpha}_n, \tilde{\beta}_n\bigr)$} with $a$ replaced by~$a/q$ gives the desired bilateral Bailey pair $(\alpha'_n, \beta'_n)$ relative to $a/q$.
\end{proof}

In addition to Lemma~\ref{lem:key1}, let us give a similarly simple lemma whose unilateral version is also due to McLaughlin in~\cite[Lemma~13.1\,(2)]{M} (see also Lovejoy \cite[Lemma 3.1]{Lo22}), and whose proof is very similar to the one of Lemma~\ref{lem:key1}.

\begin{Lemma}[key Lemma~2]\label{lem:key2}
If $(\alpha_n, \beta_n)$ is a bilateral Bailey pair relative to $a$, then $(\alpha'_n, \beta'_n)$ is a bilateral Bailey pair relative to $a/q$, where
\[%\label{eq:baileylattice2}
\alpha'_n=(1-a)\left(\frac{q^n\alpha_n}{1-aq^{2n}}-\frac{q^{n-1}\alpha_{n-1}}{1-aq^{2n-2}}\right), \qquad \beta'_n=q^n\beta_n,
\]
subject to convergence conditions on the sequences $\alpha_n$ and $\beta_n$, which make the relevant infinite series absolutely convergent.
\end{Lemma}

Using, as in the short proof of the bilateral Bailey lattice above, Lemma~\ref{lem:key2} followed by Theorem~\ref{thm:bilatbaileylemma} with $a$ replaced by $a/q$, we obtain the following new bilateral Bailey lattice, similar to Theorem~\ref{thm:bilatbaileylattice}. As far as we know, its unilateral version was also unknown until now.

\begin{Theorem}[new bilateral Bailey lattice]\label{thm:newbilatbaileylattice}
If $(\alpha_n, \beta_n)$ is a bilateral Bailey pair relative to~$a$, then $(\alpha'_n, \beta'_n)$ is a bilateral Bailey pair relative to $a/q$, where
\begin{equation*}\alpha'_n=\frac{(\rho,\sigma)_n(a/\rho\sigma)^n}{(a/\rho,a/\sigma)_n}(1-a)\left(\frac{q^n\alpha_n}{1-aq^{2n}}-\frac{q^{n-1}\alpha_{n-1}}{1-aq^{2n-2}}\right),
\end{equation*}
and
\begin{equation*}\beta'_n=\sum_{j\leq n}{(\rho,\sigma)_j(a/\rho\sigma)_{n-j}(a/\rho\sigma)^j\over (q)_{n-j}(a/\rho,a/\sigma)_n} q^j\beta_j,
\end{equation*}
subject to convergence conditions on the sequences $\alpha_n$ and $\beta_n$, which make the relevant infinite series absolutely convergent.
\end{Theorem}

Lemmas \ref{lem:key1} and \ref{lem:key2} can be generalised by adding an extra parameter $b$.
\begin{Lemma}[general lemma]\label{lem:baileylattices}
If $(\alpha_n, \beta_n)$ is a bilateral Bailey pair relative to $a$, then $(\alpha'_n, \beta'_n)$ is a bilateral Bailey pair relative to $a/q$, where
\[%\label{eq:baileylatticesalpha}
\alpha'_n=\frac{1-a}{1-b}\left(\frac{\bigl(1-bq^n\bigr) \alpha_n}{1-aq^{2n}}-\frac{q^{n-1}\bigl(aq^{n-1}-b\bigr) \alpha_{n-1}}{1-aq^{2n-2}}\right)
\qquad \text{and} \qquad \beta'_n=\frac{1-bq^n}{1-b}\beta_n.
\]
\end{Lemma}

\begin{Remark}
Lemma \ref{lem:key1} is the case $b=0$ and Lemma \ref{lem:key2} is the case $b \rightarrow \infty$ of Lemma \ref{lem:baileylattices}.
\end{Remark}

\begin{Remark}
While at first glance Lemma \ref{lem:baileylattices} seems more general than Lemmas \ref{lem:key1} and~\ref{lem:key2}, it is actually equivalent to these two lemmas taken together. Indeed, the bilateral Bailey pair in Lemma \ref{lem:baileylattices} is equal to $1/(1-b)$ times the bilateral Bailey pair of Lemma \ref{lem:key1} minus~${b/(1-b)}$ times the bilateral Bailey pair of Lemma \ref{lem:key2}. Using the fact that being a bilateral Bailey pair is stable under linear combination, Lemmas \ref{lem:key1} and \ref{lem:key2} imply Lemma \ref{lem:baileylattices}. Note that it is also possible to prove Lemma \ref{lem:baileylattices} directly with a similar method to the proof of Lemma \ref{lem:key1}.
\end{Remark}

Despite following from Lemmas \ref{lem:key1} and \ref{lem:key2}, this general Lemma \ref{lem:baileylattices} is still interesting as it provides in the unilateral case an ``inverse" to Lovejoy's Lemma 2.3 of \cite{Lo22}, which he first stated in \cite[equations (2.4) and (2.5)]{Lo04}. Actually, this result inspired our discovery of Lemma~\ref{lem:baileylattices}.


\setcounter{Theorem}{15}\setcounter{equation}{6}
% Original source lines 257--298
Moreover, from Lemma \ref{lem:baileylattices}, we deduce a very general theorem transforming bilateral Bailey pairs relative to $a$ into bilateral Bailey pairs relative to $aq^{-N}$. Recall the $M$-th elementary symmetric polynomial in $N$ variables defined for $0 \leq M \leq N$ as
\begin{equation*}e_M(X_1, \dots , X_N)= \sum_{1\leq i_1<i_2<\dots<i_M\leq N} X_{i_1} X_{i_2}\cdots X_{i_M},
\end{equation*}
and $e_M(X_1, \dots , X_N)=0$ if $M<0$ or $M>N$.

Recall also that for all $0\leq j \leq N$, the $q$-binomial coefficient is defined by
\begin{equation*}\left[{N\atop j}\right]_q=\left[{N\atop j}\right]:=\frac{(q)_N}{(q)_j(q)_{N-j}}.
\end{equation*}
We extend this definition to $j<0$ and $j>N$ by setting $\bigl[{N\atop j}\bigr]=0$, which, as $N\geq0$, is consistent with the definition of $q$-Pochhammer symbols with negative indices given above.
The general theorem can be stated as follows.

\begin{Theorem}[new bilateral Bailey lattice in higher dimension]\label{thm:multibaileylattice}
Let $(\alpha_n, \beta_n)$ be a bilateral Bailey pair relative to $a$. For all $N\geq 1$, define the pair \smash{$\bigl(\alpha^{(N)}_n, \beta^{(N)}_n\bigr)$} by
\begin{equation}
\alpha^{(N)}_n=\frac{\bigl(1-aq^{2n-N}\bigr)\bigl(aq^{1-N}\bigr)_N}{(1-b_1)\cdots(1-b_N)}\sum_{j \in \Z}(-1)^j\frac{q^{jn-j(j+1)/2}f_{N,j,n}(b_1,\ldots,b_N)}{\bigl(aq^{2n-N-j}\bigr)_{N+1}} \alpha_{n-j},\label{eq:multibaileylatticealpha}
\end{equation}
where
\begin{align}
f_{N,j,n}(b_1,\ldots,b_N) :={}& \sum_{M \in \Z} \sum_{u \in \Z} a^{u}q^{(M-j+u)(n-j+u)+u(n-N)} \begin{bmatrix}
M\\
j-u
\end{bmatrix} \begin{bmatrix}
N-M\\
u
\end{bmatrix}\nonumber\\
&\times (-1)^Me_M(b_1, \dots , b_N),\label{eq:pascalab}
\end{align}
and
\begin{equation}
\beta^{(N)}_n=\left(\prod_{i=1}^N \frac{1-b_iq^n}{1-b_i}\right)\beta_n.\label{eq:multibaileylatticebeta}
\end{equation}
Then, \smash{$\bigl(\alpha^{(N)}_n, \beta^{(N)}_n\bigr)$} is a bilateral Bailey pair relative to $aq^{-N}$.
\end{Theorem}

\begin{Remark}
When $j<0$ or $j > N$, we have $f_{N,j,n}(b_1,\ldots,b_N) =0$ because of the $q$-binomial coefficients in~\eqref{eq:pascalab}. Therefore, the sum in \eqref{eq:multibaileylatticealpha} is actually finite.
\end{Remark}

\begin{Remark}
Note that the sums over $M$ and $u$ in \eqref{eq:pascalab} could equivalently be taken from~$0$ to $N$ and from~$0$ to $j$, respectively. Indeed the $q$-binomial coefficients or $e_M(b_1,\dots ,b_N)$ naturally cancel outside of these ranges. However, the expression with infinite sums makes future calculations easier to write.
\end{Remark}


\setcounter{section}{3}\setcounter{subsection}{1}\setcounter{Theorem}{6}\setcounter{equation}{0}
% Original source lines 963--1257
\subsection[Proof of Theorem~\ref{thm:multibaileylattice}]{Proof of Theorem~\ref{thm:multibaileylattice}}

Now we can turn to the proof of Theorem~\ref{thm:multibaileylattice}.
Recall the classical $q$-analogues of Pascal's triangle
\begin{equation}\label{eq:pascal1}
\left[{N+1\atop j}\right] = q^j\left[{N\atop j}\right]+\left[{N\atop j-1}\right],
\end{equation}
and
\begin{equation} \label{eq:pascal2}
\left[{N+1\atop j}\right]=\left[{N\atop j}\right]+ q^{N+1-j}\left[{N\atop j-1}\right],
\end{equation}
for all integers $N$, $j$ with $N\geq 0$.

\subsubsection[Recurrence relation for f\_\{N,j,n\}(b\_1,...,b\_N)]{Recurrence relation for $\boldsymbol{f_{N,j,n}(b_1,\ldots,b_N)}$}
We first prove the following recurrence relation, which plays a central role in our proof.
\begin{Proposition}[Recurrence relation]\label{prop:rec} For all $0\leq j\leq N+1$,
\begin{gather*}
\bigl(1-aq^{2n-1-N}\bigr)f_{N+1,j,n}(b_1,\ldots,b_{N+1})\\
\qquad =(1-b_{N+1}q^n)\bigl(1-aq^{2n-1-N-j}\bigr)f_{N,j,n}(b_1,\ldots,b_N)\\
\phantom{\qquad =}{}+\bigl(aq^{n-1-N}-b_{N+1}\bigr)\bigl(1-aq^{2n-j}\bigr)f_{N,j-1,n-1}(b_1,\ldots,b_N).
\end{gather*}
\end{Proposition}

We start by giving two technical lemmas which are extensions of \eqref{eq:pascal1} and \eqref{eq:pascal2}. Their proofs are straightforward verifications and are therefore omitted here.


\begin{Lemma}
\label{lem:tech1}
For all $j,u \in \Z$ and $0\leq M\leq N$,
\begin{gather*}
\begin{bmatrix}
M\\
j-u
\end{bmatrix} \begin{bmatrix}
N+1-M\\
u
\end{bmatrix}-q^{2j-2u-M+1}\begin{bmatrix}
M\\
j-u+1
\end{bmatrix} \begin{bmatrix}
N+1-M\\
u-1
\end{bmatrix}\\
\qquad=q^{u}\begin{bmatrix}
M\\
j-u
\end{bmatrix} \begin{bmatrix}
N-M\\
u
\end{bmatrix}
+q^{j-u-M}\begin{bmatrix}
N-M\\
u-1
\end{bmatrix}\left(\begin{bmatrix}
M\\
j-u
\end{bmatrix}-\begin{bmatrix}
M\\
j-u+1
\end{bmatrix}\right)\\
\phantom{\qquad=}{} -q^{2j-3u-2M+3+N}\begin{bmatrix}
M\\
j-u+1
\end{bmatrix} \begin{bmatrix}
N-M\\
u-2
\end{bmatrix} .
\end{gather*}
\end{Lemma}

Note that when $M=0$ and $u=j$, Lemma~\ref{lem:tech1} reduces to \eqref{eq:pascal1}.

\begin{Lemma}
\label{lem:tech2}
For all $j,u \in \Z$ and $0\leq M\leq N$,
\begin{gather*}
\begin{bmatrix}
M+1\\
j-u
\end{bmatrix} \begin{bmatrix}
N-M\\
u
\end{bmatrix}-q^{2j-2u-M}\begin{bmatrix}
M+1\\
j-u+1
\end{bmatrix} \begin{bmatrix}
N-M\\
u-1
\end{bmatrix}\\
\qquad
=q^j \begin{bmatrix}
M\\
j-u
\end{bmatrix} \begin{bmatrix}
N-M\\
u
\end{bmatrix}
- q^{2j-2u-M}\begin{bmatrix}
M\\
j-u+1
\end{bmatrix} \begin{bmatrix}
N-M\\
u-1
\end{bmatrix}+\begin{bmatrix}
M\\
j-u-1
\end{bmatrix} \begin{bmatrix}
N-M\\
u
\end{bmatrix}\\
\phantom{\qquad=}{}- q^{N+j+1-2u-M} \begin{bmatrix}
M\\
j-u
\end{bmatrix} \begin{bmatrix}
N-M\\
u-1
\end{bmatrix}.
\end{gather*}
\end{Lemma}

Note that when $M=1$ and $u=j$, Lemma~\ref{lem:tech2} reduces to a combination of \eqref{eq:pascal1} and \eqref{eq:pascal2}.

We can now prove our recurrence relation.
\begin{proof}[Proof of Proposition \ref{prop:rec}]
Using the homogeneity of $e_M$, recall from Theorem~\ref{thm:multibaileylattice} that
\begin{gather*}
f_{N,j,n}(b_1,\ldots,b_N)\\
\qquad= \sum_{M \in \Z} \sum_{u \in \Z} a^{u}q^{(M-j+u)(n-j+u)+u(n-N)} \begin{bmatrix}
M\\
j-u
\end{bmatrix} \begin{bmatrix}
N-M\\
u
\end{bmatrix} e_M(-b_1, \dots , -b_N).
\end{gather*}
Thus
\begin{gather}
\bigl(1-aq^{2n-1-N}\bigr)f_{N+1,j,n}(b_1,\ldots,b_{N+1})\nonumber \\
\qquad= \sum_{M \in \Z} \sum_{u\in \mathbb{Z}} a^{u}q^{(M-j+u)(n-j+u)+u(n-N-1)}\nonumber\\
\phantom{\qquad=}{}\times\bigg( \begin{bmatrix}
M\\
j-u
\end{bmatrix} \begin{bmatrix}
N+1-M\\
u
\end{bmatrix}
-q^{2j-2u-M+1}\begin{bmatrix}
M\\
j-u+1
\end{bmatrix} \begin{bmatrix}
N+1-M\\
u-1
\end{bmatrix} \bigg)\nonumber\\
\phantom{\qquad=}{}\times e_M(-b_1, \dots , -b_{N+1}).\label{eq:lastline}
\end{gather}
Writing
\begin{equation*}e_M(-b_1, \dots , -b_{N+1})=e_M(-b_1, \dots , -b_{N})-b_{N+1}e_{M-1}(-b_1, \dots , -b_{N}),
\end{equation*}
it is enough to show that coefficients of $b_{N+1}^0a^ue_M(-b_1, \dots , -b_{N})$ and $b_{N+1}^1a^ue_M(-b_1, \dots , -b_{N})$ coincide on both sides of Proposition~\ref{prop:rec}.

First, we derive from~\eqref{eq:lastline} that the coefficient of $b_{N+1}^0a^ue_M(-b_1, \dots , -b_{N})$ on the left-hand side of Proposition~\ref{prop:rec} is equal to \smash{$q^{(M-j+u)(n-j+u)+u(n-N-1)}$} times
\begin{equation}\label{eq:coeffb0l}
\begin{bmatrix}
M\\
j-u
\end{bmatrix} \begin{bmatrix}
N+1-M\\
u
\end{bmatrix}-q^{2j-2u-M+1}\begin{bmatrix}
M\\
j-u+1
\end{bmatrix} \begin{bmatrix}
N+1-M\\
u-1
\end{bmatrix}.
\end{equation}

Now the coefficient of $b_{N+1}^0a^ue_M(-b_1, \dots , -b_{N})$ on the right-hand side of Proposition~\ref{prop:rec} is equal to \smash{$q^{(M-j+u)(n-j+u)+u(n-N-1)}$} times
\begin{gather}
q^u \begin{bmatrix}
M\\
j-u
\end{bmatrix} \begin{bmatrix}
N-M\\
u
\end{bmatrix}-q^{j-u-M}\begin{bmatrix}
M\\
j-u+1
\end{bmatrix} \begin{bmatrix}
N-M\\
u-1
\end{bmatrix}+q^{j-u-M}\begin{bmatrix}
M\\
j-u
\end{bmatrix} \begin{bmatrix}
N-M\\
u-1
\end{bmatrix}\nonumber\\
\qquad{}-q^{N+2j-2M-3u+3} \begin{bmatrix}
M\\
j-u+1
\end{bmatrix} \begin{bmatrix}
N-M\\
u-2
\end{bmatrix}.\label{eq:coeffb0r}
\end{gather}
Note finally that~\eqref{eq:coeffb0l} (resp.~\eqref{eq:coeffb0r}) is the left-hand (resp.\ right-hand) side of Lemma~\ref{lem:tech1}, so we are done regarding this first coefficient.

Similarly, the coefficient of $b_{N+1}^1a^ue_M(-b_1, \dots , \!-b_{N})$ in~\eqref{eq:lastline} is \smash{$-q^{(M+1-j+u)(n-j+u)+u(n-N-1)}$} times
\begin{equation}\label{eq:coeffb1l}
\begin{bmatrix}
M+1\\
j-u
\end{bmatrix} \begin{bmatrix}
N-M\\
u
\end{bmatrix}-q^{2j-2u-M}\begin{bmatrix}
M+1\\
j-u+1
\end{bmatrix} \begin{bmatrix}
N-M\\
u-1
\end{bmatrix},
\end{equation}

and the coefficient of $b_{N+1}^1a^ue_M(-b_1, \dots , -b_{N})$ on the right-hand side of Proposition~\ref{prop:rec} is equal to \smash{$-q^{(M+1-j+u)(n-j+u)+u(n-N-1)}$} times
\begin{gather}
q^j \begin{bmatrix}
M\\
j-u
\end{bmatrix} \begin{bmatrix}
N-M\\
u
\end{bmatrix}-q^{2j-2u-M}\begin{bmatrix}
M\\
j-u+1
\end{bmatrix} \begin{bmatrix}
N-M\\
u-1
\end{bmatrix}+\begin{bmatrix}
M\\
j-u-1
\end{bmatrix} \begin{bmatrix}
N-M\\
u
\end{bmatrix}\nonumber\\
\qquad{}-q^{N+j+1-2u-M} \begin{bmatrix}
M\\
j-u
\end{bmatrix} \begin{bmatrix}
N-M\\
u-1
\end{bmatrix}.\label{eq:coeffb1r}
\end{gather}

As~\eqref{eq:coeffb1l} (resp.~\eqref{eq:coeffb1r}) is the left-hand (resp.\ right-hand) side of Lemma~\ref{lem:tech2}, this ends the proof of Proposition~\ref{prop:rec}.
\end{proof}

\subsubsection[Proof of Theorem \ref{thm:multibaileylattice}]{Proof of Theorem \ref{thm:multibaileylattice}}

Now we use Proposition \ref{prop:rec} and Lemma \ref{lem:baileylattices} to prove Theorem \ref{thm:multibaileylattice}.

We proceed by induction on $N$. For $N=0$, by \eqref{eq:multibaileylatticebeta}, \smash{$\beta^{(0)}_n=\beta_n$} and by \eqref{eq:multibaileylatticealpha},
\begin{equation*}
\alpha^{(0)}_n = \bigl(1-aq^{2n}\bigr)\frac{\alpha_{n}}{1-aq^{2n}}=\alpha_n.
\end{equation*}
Thus $\bigl(\alpha^{(0)}_n, \beta^{(0)}_n\bigr)$ is a bilateral Bailey pair relative to $a$.

Now, assume that for some integer $N\geq 0$, \smash{$\bigl(\alpha^{(N)}_n,\beta^{(N)}_n\bigr)$} is a bilateral Bailey pair relative to~$aq^{-N}$ and show that \smash{$\bigl(\alpha^{(N+1)}_n,\beta^{(N+1)}_n\bigr)$} is a bilateral Bailey pair relative to $aq^{-N-1}$. By Lemma~\ref{lem:baileylattices} with $b = b_{N+1}$, $(\alpha'_n,\beta'_n)$ is a bilateral Bailey pair relative to $aq^{-N-1}$, where
\begin{equation*}
\alpha'_n=\frac{1-aq^{-N}}{1-b_{N+1}}\left(\frac{(1-b_{N+1}q^n) \alpha_n^{(N)}}{1-aq^{2n-N}}-\frac{q^{n-1}\bigl(aq^{n-N-1}-b_{N+1}\bigr) \alpha_{n-1}^{(N)}}{1-aq^{2n-N-2}}\right),
\end{equation*}
and
\begin{equation*}\beta'_n=\frac{1-b_{N+1}q^n}{1-b_{N+1}}\beta_n^{(N)}.
\end{equation*}

Now let us show that $(\alpha'_n,\beta'_n) = \bigl(\alpha^{(N+1)}_n,\beta^{(N+1)}_n\bigr)$.
It is clear that $\beta'_n=\beta^{(N+1)}_n$, and for $\alpha'_n$, by~\eqref{eq:multibaileylatticealpha} we have
\begin{align*}
\alpha'_n
={}&\frac{1-aq^{-N}}{1-b_{N+1}}\left(\frac{(1-b_{N+1}q^n) \alpha_n^{(N)}}{1-aq^{2n-N}}-\frac{q^{n-1}\bigl(aq^{n-N-1}-b_{N+1}\bigr) \alpha_{n-1}^{(N)}}{1-aq^{2n-N-2}}\right)\\
={}& \frac{\bigl(aq^{-N}\bigr)_{N+1}}{(1-b_1)\cdots(1-b_{N+1})} \bigg(\sum_{j \in \Z}(-1)^j\frac{q^{jn-j(j+1)/2}(1-b_{N+1}q^n)f_{N,j,n}(b_1,\ldots,b_N)}{\bigl(aq^{2n-N-j}\bigr)_{N+1}} \alpha_{n-j} \\
& +\sum_{j \in \Z}(-1)^{j+1}\frac{q^{(j+1)(n-1)-j(j+1)/2}\bigl(aq^{n-1-N}-b_{N+1}\bigr)f_{N,j,n-1}(b_1,\ldots,b_N)}{\bigl(aq^{2n-2-N-j}\bigr)_{N+1}} \alpha_{n-1-j} \bigg)\\
={}& \frac{\bigl(aq^{-N}\bigr)_{N+1}}{(1-b_1)\cdots(1-b_{N+1})} \sum_{j \in \Z}(-1)^j\frac{q^{jn-j(j+1)/2}}{(aq^{2n-N-j-1})_{N+2}} \big( (1-b_{N+1}q^n)\bigl(1-aq^{2n-N-j-1}\bigr)
\\
& \times f_{N,j,n}(b_1,\ldots,b_N)+ \bigl(aq^{n-1-N}-b_{N+1}\bigr)\bigl(1-aq^{2n-j}\bigr)f_{N,j-1,n-1}(b_1,\ldots,b_N) \big) \alpha_{n-j}.
\end{align*}
Now by Proposition \ref{prop:rec}, this equals
\begin{align*}
\alpha'_n
={}& \frac{\bigl(aq^{-N}\bigr)_{N+1}}{(1-b_1)\cdots(1-b_{N+1})} \sum_{j \in \Z}(-1)^j\frac{q^{jn-j(j+1)/2}}{\bigl(aq^{2n-N-j-1}\bigr)_{N+2}} \bigl(1-aq^{2n-1-N}\bigr)\\
&\times f_{N+1,j,n}(b_1,\ldots,b_{N+1}) \alpha_{n-j}
= \alpha_n^{(N+1)}.
\end{align*}
Thus the pair \smash{$\bigl(\alpha^{(N+1)}_n, \beta^{(N+1)}_n\bigr)$} is indeed a Bailey pair relative to $aq^{-N-1}$.
\begin{thebibliography}{99}
\bibitem[4]{A1}
Andrews G.E., Multiple series {R}ogers--{R}amanujan type identities,
 \href{https://doi.org/10.2140/pjm.1984.114.267}{\textit{Pacific~J. Math.}} \textbf{114} (1984), 267--283.

\bibitem[5]{A2}
Andrews G.E., {$q$}-series: their development and application in analysis,
 number theory, combinatorics, physics, and computer algebra, \textit{CBMS
 Reg. Conf. Ser. Math.}, Vol.~66, American Mathematical Society, Providence,
 RI, 1986.

\bibitem[6]{AAR}
Andrews G.E., Askey R., Roy R., Special functions, \textit{Encyclopedia Math.
 Appl.}, Vol.~71, \href{https://doi.org/10.1017/CBO9781107325937}{Cambridge University Press}, Cambridge, 1999.

\bibitem[8]{Ba}
Bailey W.N., Identities of the {R}ogers--{R}amanujan type, \href{https://doi.org/10.1112/plms/s2-50.1.1}{\textit{Proc. London
 Math. Soc.}} \textbf{50} (1948), 1--10.

\bibitem[9]{BMS}
Berkovich A., McCoy B.M., Schilling A., {$N=2$}~supersymmetry and {B}ailey
 pairs, \href{https://doi.org/10.1016/S0378-4371(96)00054-4}{\textit{Phys.~A}} \textbf{228} (1996), 33--62, \href{https://arxiv.org/abs/hep-th/9512182}{arXiv:hep-th/9512182}.

\bibitem[16]{GR}
Gasper G., Rahman M., Basic hypergeometric series, 2nd ed., \textit{Encyclopedia Math.
 Appl.}, Vol.~96, \href{https://doi.org/10.1017/CBO9780511526251}{Cambridge University Press}, Cambridge, 2004.

\bibitem[21]{J}
Jouhet F., Shifted versions of the {B}ailey and well-poised {B}ailey lemmas,
 \href{https://doi.org/10.1007/s11139-010-9282-x}{\textit{Ramanujan~J.}} \textbf{23} (2010), 315--333, \href{https://arxiv.org/abs/0906.1870}{arXiv:0906.1870}.

\bibitem[24]{Lo04}
Lovejoy J., A {B}ailey lattice, \href{https://doi.org/10.1090/S0002-9939-03-07247-2}{\textit{Proc. Amer. Math. Soc.}} \textbf{132}
 (2004), 1507--1516.

\bibitem[25]{Lo22}
Lovejoy J., Bailey pairs and indefinite quadratic forms,~{II}. {F}alse
 indefinite theta functions, \href{https://doi.org/10.1007/s40993-022-00324-x}{\textit{Res. Number Theory}} \textbf{8} (2022),
 24, 16~pages.

\bibitem[26]{M}
McLaughlin J., Topics and methods in {$q$}-series, \textit{Monogr. Number
 Theory}, Vol.~8, World Scientific Publishing, Hackensack, NJ, 2018.

\end{thebibliography}
\end{document}
